\subsection{5G LENA NR Module}
5G-LENA NR \cite{5glena} is an extension of ns-3 developed by CTTC for simulating 3GPP compliant 5G NR cellular networks across both  sub-6 Hz and mmWave frequency ranges (0.5 to 100 GHz).
The module implements  5G NR core network components, protocol stack including the Physical (PHY), Medium Access Control (MAC), Radio Link Control (RLC), and Packet Data Convergence Protocol (PDCP) layers.
The MAC layer in 5G-LENA supports flexible schedulers (e.g., Proportional Fair, Round Robin) for both TDMA and OFDMA-based access, while the PHY layer handles advanced MIMO models and MCS (Modulation and Coding Scheme) adaptation.

While 5G-LENA NR provides the protocol stacks for 5G NR, it still rely on the standard ns-3 propagation models which cause the same limitations previously discussed.
Specifically, it lacks the site-specific channel information related to environment that it is deployed in.
Such information are important to model different research scenarios such as ISAC in 6G network systems.
This gap in physical-layer realism is the final motivator for us to integrate a ray-tracing engine into 5G-LENA NR.